\section*{Chapter \ref{ch:g11:func} --- Functions and Variations}

\begin{solution}{exo:g11:func:1}
$f$: need $3 - x \geq 0$, \omterm{def:g11:func:function}{domain} $\intoc{-\infty}{3}$.

$g$: need $x^2 - 4 \neq 0$, i.e.\ $x \neq \pm 2$: \omterm{def:g11:func:function}{domain}
$\R \setminus \{-2, 2\}$.

$h$: need $x \geq 0$ and $x \neq 5$: \omterm{def:g11:func:function}{domain}
$\intco{0}{5} \cup \intoo{5}{+\infty}$.
\end{solution}

\begin{solution}{exo:g11:func:2}
$\sqrt{11} < \sqrt{13}$ since $\sqrt{}$ is \omterm{def:g11:func:monotone}{increasing} and $11 < 13$.

$(-2.1)^2 > (-2.05)^2$ since $x^2$ is \omterm{def:g10:functions:variations}{decreasing} on
$\intoc{-\infty}{0}$ and $-2.1 < -2.05$.

$\frac{1}{0.9} > \frac{1}{1.1}$ since $\frac1x$ is \omterm{def:g10:functions:variations}{decreasing} on
$\intoo{0}{+\infty}$ and $0.9 < 1.1$.

$(-1.01)^3 < (-0.99)^3$ since $x^3$ is \omterm{def:g11:func:monotone}{increasing} and $-1.01 < -0.99$.
\end{solution}

\begin{solution}{exo:g11:func:3}
$f(-x) = 3(-x)^4 - (-x)^2 = 3x^4 - x^2 = f(x)$: \omterm{def:g11:func:parity}{even}.

$g(-x) = -x^3 - x = -g(x)$: \omterm{def:g11:func:parity}{odd}.

$h(1) = 2$ and $h(-1) = 1 - 1 = 0$, so $h(-1)$ is neither $h(1)$ nor
$-h(1)$: neither \omterm{def:g11:func:parity}{even} nor \omterm{def:g11:func:parity}{odd}.

$k(-x) = \frac{-x}{x^2+1} = -k(x)$: \omterm{def:g11:func:parity}{odd}.
\end{solution}

\begin{solution}{exo:g11:func:4}
Read the \omterm{def:g10:functions:table}{variation table}. On $\intcc{-3}{1}$, $f$ climbs from $-2$ to $4$
and crosses the value $1$ exactly once; on $\intcc{1}{5}$ it goes down
from $4$ to $0$ and crosses $1$ once more: the \omterm{def:g10:algebra:equation}{equation} $f(x) = 1$ has
$2$ solutions. The \omterm{def:g10:functions:extrema}{maximum} of $f$ is $f(1) = 4 < 5$, so $f(x) = 5$ has no
solution.
\end{solution}

\begin{solution}{exo:g11:func:5}
Let $2 \leq u < v$. Then
\[
f(v) - f(u) = v^2 - 4v - u^2 + 4u = (v-u)(v+u) - 4(v-u)
= (v - u)(u + v - 4).
\]
Both factors are positive: $v - u > 0$, and $u + v > 2 + 2 = 4$ since
$u \geq 2$ and $v > 2$. Hence $f(v) > f(u)$: $f$ is strictly \omterm{def:g11:func:monotone}{increasing}
on $\intco{2}{+\infty}$.
\end{solution}

\begin{solution}{exo:g11:func:6}
$f = \sqrt{u}$ with $u(x) = 5 - x$ \omterm{def:g10:functions:variations}{decreasing} and nonnegative on
$\intoc{-\infty}{5}$: $f$ is \omterm{def:g10:functions:variations}{decreasing} there
(\cref{prop:g11:func:transforms}).

$g = \frac1u$ with $u(x) = x^2 + 1$ \omterm{def:g11:func:monotone}{increasing} and positive on
$\intco{0}{+\infty}$: $g$ is \omterm{def:g10:functions:variations}{decreasing} there.

$h = 3 - 2\sqrt x$: $\sqrt x$ is \omterm{def:g11:func:monotone}{increasing}, multiplying by $-2 < 0$
reverses, adding $3$ changes nothing: $h$ is \omterm{def:g10:functions:variations}{decreasing} on
$\intco{0}{+\infty}$.
\end{solution}

\begin{solution}{exo:g11:func:7}
\emph{1.} $f$ is in \omterm{thm:g11:quad:canonical}{canonical form}: \omterm{def:g10:functions:variations}{decreasing} on $\intoc{-\infty}{2}$,
\omterm{def:g11:func:monotone}{increasing} on $\intco{2}{+\infty}$, \omterm{def:g10:functions:extrema}{minimum} $f(2) = -3$.

\emph{2.} $g = \frac{1}{u}$ with $u(x) = (x-2)^2 + 1 = f(x) + 4$. The
shift by $4$ keeps the variations of $f$, and $u \geq 1 > 0$. Taking the
reciprocal reverses them: $g$ is \omterm{def:g11:func:monotone}{increasing} on $\intoc{-\infty}{2}$ and
\omterm{def:g10:functions:variations}{decreasing} on $\intco{2}{+\infty}$, with \omterm{def:g10:functions:extrema}{maximum}
$g(2) = \frac{1}{1} = 1$.
\end{solution}

\begin{solution}{exo:g11:func:8}
Any hill-shaped curve through the three points works. Then:
$f + 2$ passes through $(0,3)$, $(2,5)$, $(4,3)$, same variations
(vertical shift); $-f$ passes through $(0,-1)$, $(2,-3)$, $(4,-1)$ with
\emph{reversed} variations (the peak becomes a valley); $2f$ passes
through $(0,2)$, $(2,6)$, $(4,2)$, same variations but amplified
vertically.
\end{solution}

\begin{solution}{exo:g11:func:9}
$2 - \frac{1}{x+1} = \frac{2(x+1) - 1}{x+1} = \frac{2x+1}{x+1} = f(x)$.
On $\intoo{-1}{+\infty}$, $u(x) = x + 1$ is \omterm{def:g11:func:monotone}{increasing} and positive, so
$\frac1u$ is \omterm{def:g10:functions:variations}{decreasing}, $-\frac1u$ is \omterm{def:g11:func:monotone}{increasing}, and adding $2$ changes
nothing: $f$ is strictly \omterm{def:g11:func:monotone}{increasing} on $\intoo{-1}{+\infty}$.
\end{solution}

\begin{solution}{exo:g11:func:10}
\emph{1. True.} If $u < v$, then
$(f+g)(v) - (f+g)(u) = \bigl(f(v)-f(u)\bigr) + \bigl(g(v)-g(u)\bigr)$ is
a sum of two nonnegative numbers.

\emph{2. False.} $f(x) = g(x) = x$ are \omterm{def:g11:func:monotone}{increasing} on $\R$, but their
product $x^2$ is not \omterm{def:g11:func:monotone}{increasing} on $\R$ (it decreases on
$\intoc{-\infty}{0}$). (The statement becomes true if both \omterm{def:g11:func:function}{functions} are
also nonnegative.)

\emph{3. False.} Same counterexample: $f(x) = x$ is \omterm{def:g11:func:monotone}{increasing} on $\R$
but $f^2(x) = x^2$ is not.
\end{solution}

\begin{solution}{exo:g11:func:11}
\emph{1.} For $0 < u < v$:
\[
f(v) - f(u) = (v - u) + \frac1v - \frac1u
= (v - u) + \frac{u - v}{uv}
= (v - u)\left(1 - \frac{1}{uv}\right)
= (v - u)\,\frac{uv - 1}{uv}.
\]

\emph{2.} Always $v - u > 0$ and $uv > 0$. If $0 < u < v \leq 1$, then
$uv < 1$, so $uv - 1 < 0$ and $f(v) < f(u)$: $f$ is strictly \omterm{def:g10:functions:variations}{decreasing}
on $\intoc{0}{1}$. If $1 \leq u < v$, then $uv > 1$ and $f(v) > f(u)$:
$f$ is strictly \omterm{def:g11:func:monotone}{increasing} on $\intco{1}{+\infty}$.

\emph{3.} $f$ decreases before $1$ and increases after, so
$f(x) \geq f(1) = 2$ for all $x > 0$, with equality exactly at $x = 1$.
(This inequality reappears in many guises; compare
\cref{ex:g11:quad:canonical} and the tangent-line inequalities of
\cref{ch:g12:deriv}.)
\end{solution}
